% Generated draft: requires mathematical and source review.
\documentclass[11pt]{article}
\usepackage[margin=25mm]{geometry}
\usepackage{fontspec}
\setmainfont{Latin Modern Roman}
\usepackage{amsmath,amssymb,mathtools,unicode-math}
\setmathfont{Latin Modern Math}
\usepackage{enumitem,needspace,xurl,hyperref}
\setlength{\parindent}{0pt}
\setlength{\parskip}{5pt}
\setlength{\emergencystretch}{3em}
\newcommand{\N}{\mathbb{N}}
\newcommand{\Z}{\mathbb{Z}}
\newcommand{\Q}{\mathbb{Q}}
\newcommand{\R}{\mathbb{R}}
\newcommand{\C}{\mathbb{C}}
\newcommand{\F}{\mathbb{F}}
\newcommand{\A}{\mathbb{A}}
\newcommand{\PP}{\mathbb{P}}
\newcommand{\E}{\mathbb{E}}
\newcommand{\eps}{\varepsilon}
\newcommand{\dd}{\,\mathrm{d}}
\begin{document}
\section*{UnsolvedMath ID: 30000312}
\textbf{Source problem code:} catalogue entry
\subsection{Definitions and mathematical statement}
Let $\mathcal C$ be the class in the restored statement: nonzero simple, separable, nuclear, stable $C^*$-algebras $A$ of real rank zero. Simplicity means absence of nonzero proper closed two-sided ideals; stability means $A\cong A\otimes\mathcal K$, where $\mathcal K$ denotes the compact operators on a separable infinite-dimensional Hilbert space. Real rank zero means that finite-spectrum self-adjoint elements are norm dense in the self-adjoint part of the unitization, using the usual nonunital convention. Nuclearity is the usual uniqueness of the $C^*$-tensor norm with arbitrary $C^*$-algebras.

For projections, Murray--von Neumann equivalence is implemented by a partial isometry. A projection $p$ is finite if $q\le p$ and $q\sim p$ imply $q=p$; otherwise it is infinite. In this real-rank-zero setting the stable-finiteness question is expressed by finiteness of all projections in matrix algebras over $A$. For a simple algebra the usual projection-based pure-infiniteness condition is that every nonzero hereditary $C^*$-subalgebra contain an infinite projection. The source's exact nonunital conventions should be checked.

Let $K_0(A)$ be the usual Grothendieck $K$-group of projections, using unitization for a nonunital algebra, let $K_0(A)^+$ be its cone of classes represented by projections over $A$, and let $K_1(A)$ be the usual topological $K_1$ group. Write
\[
 I(A)=(K_0(A),K_0(A)^+,K_1(A)).
\]
The three questions are:
\begin{enumerate}
\item For every $A\in\mathcal C$, is $A$ stably finite or purely infinite?
\item For every $A\in\mathcal C$, is $K_0(A)$ weakly unperforated? With the common strict-positivity convention, this asks whether
\[
 nx\in K_0(A)^+\setminus\{0\},\quad n\ge1
 \quad\Longrightarrow\quad x\in K_0(A)^+\setminus\{0\}.
\]
The intended convention must be verified, particularly if the positive cone is not proper.
\item Is $I$ a complete invariant on $\mathcal C$? At the object-classification level this asks whether, for every $A,B\in\mathcal C$, an ordered-group isomorphism $K_0(A)\cong K_0(B)$ carrying positive cone onto positive cone, together with a group isomorphism $K_1(A)\cong K_1(B)$, implies $A\cong B$ as $C^*$-algebras. No unit class, trace data, or lifting of each prescribed invariant isomorphism is added.
\end{enumerate}

The older original wording appears to attach the separable, nuclear, stable restriction explicitly only to the final question, whereas the restored statement puts all three questions in $\mathcal C$. The supplied verification says the ambient class was made explicit from the surviving extract. The scope of the first two questions therefore remains a source-review issue; this draft displays the restored formulation without erasing that discrepancy.
\subsection{Short English statement}
For the restored class of simple separable nuclear stable real-rank-zero $C^*$-algebras, the source asks for a finite/purely-infinite dichotomy, weak unperforation of ordered $K_0$, and classification by ordered $K_0$ together with $K_1$. The original wording leaves the scope of the first two questions less explicit.
\subsection{Sources}
\begin{itemize}
\item Local UnsolvedMath registry, permanent ID 30000312, code catalogue entry, ``Classification of Simple Real-Rank-Zero $C^*$-Algebras.'' Proposed year 2005; proposer unspecified.
\item \emph{$C^*$-Algebren} (2006), Oberwolfach Reports, pp.\ 2305--2374, DOI 10.4171/owr/2005/41; \url{https://doi.org/10.4171/owr/2005/41}. The displayed report year 2006 differs from the DOI and proposed year 2005.
\item Statement status corrected\_verified; the supplied verification says the ambient class and all three questions were made explicit from the surviving extract.
\item Literature triage checked 2026-08-21: open. It says the three questions are listed as open in the source and no general result closes all of them. This bundled dated assessment is not a verified status for each individual question. Created 2026-08-25; updated 2026-09-16. Status and cited claims require independent verification.
\item Registry curation by Ulam AI / UnsolvedMath Contributors, CC BY 4.0; the underlying report retains its supplied provenance. Unreviewed statement-only draft.
\end{itemize}
\end{document}
